% ============================================================
% content/questions.tex —— 2021 湖南卷物理 · 试题
% 来源：2021/2021_物理真题_含解析.md（来源站 Word 版转录）
% 题号保留官方编号；图形见 figures/（TikZ 重绘）与 photos/（非矢量例外）
% 注：官方卷第 15、16 题为选考题（3-3/3-4 二选一），本工程转录第 15 题（选修 3-3）
% ============================================================

\gkpart{一、选择题}

\textbf{（一）单项选择题（本题共 6 小题，每小题 4 分，共 24 分。在每小题给出的四个选项中，只有一项是符合题目要求的。）}

\begin{question}
\gkti{1}核废料具有很强的放射性，需要妥善处理。下列说法正确的是\gkblank

\begin{gkopts}(2)
  \task A．放射性元素经过两个完整的半衰期后，将完全衰变殆尽
  \task B．原子核衰变时电荷数守恒，质量数不守恒
  \task C．改变压力、温度或浓度，将改变放射性元素的半衰期
  \task D．过量放射性辐射对人体组织有破坏作用，但辐射强度在安全剂量内则没有伤害
\end{gkopts}
\end{question}

\begin{question}
\gkti{2}物体的运动状态可用位置~$x$~和动量~$p$~描述，称为相，对应~$p$--$x$~图像中的一个点。物体运动状态的变化可用~$p$--$x$~图像中的一条曲线来描述，称为相轨迹。假如一质点沿~$x$~轴正方向做初速度为零的匀加速直线运动，则对应的相轨迹可能是\gkblank

\gkfig{figures/q2-phase}
\end{question}

\begin{question}
\gkti{3}``复兴号''动车组用多节车厢提供动力，从而达到提速的目的。总质量为~$m$~的动车组在平直的轨道上行驶。该动车组有四节动力车厢，每节车厢发动机的额定功率均为~$P$，若动车组所受的阻力与其速率成正比（$F_{\text{阻}}=kv$，$k$~为常量），动车组能达到的最大速度为~$v_{m}$。下列说法正确的是\gkblank

\begin{gkopts}(1)
  \task A．动车组在匀加速启动过程中，牵引力恒定不变
  \task B．若四节动力车厢输出功率均为额定值，则动车组从静止开始做匀加速运动
  \task C．若四节动力车厢输出的总功率为 $2.25P$，则动车组匀速行驶的速度为 $\dfrac{3}{4}v_{m}$
  \task D．若四节动力车厢输出功率均为额定值，动车组从静止启动，经过时间~$t$~达到最大速度~$v_{m}$，则这一过程中该动车组克服阻力做的功为 $\dfrac{1}{2}mv_{m}^{2}-Pt$
\end{gkopts}
\end{question}

\begin{question}
\gkti{4}如图所示，在~$(a,0)$~位置放置电荷量为~$q$~的正点电荷，在~$(0,a)$~位置放置电荷量为~$q$~的负点电荷，在距~$P(a,a)$~为~$\sqrt{2}\,a$~的某点处放置正点电荷~$Q$，使得~$P$~点的电场强度为零。则~$Q$~的位置及电荷量分别为\gkblank

\gkfig{figures/q4-charges}

\begin{gkopts}(2)
  \task A．$(0,2a)$，$\sqrt{2}\,q$
  \task B．$(0,2a)$，$2\sqrt{2}\,q$
  \task C．$(2a,0)$，$\sqrt{2}\,q$
  \task D．$(2a,0)$，$2\sqrt{2}\,q$
\end{gkopts}
\end{question}

\begin{question}
\gkti{5}质量为~$m'$~的凹槽静止在水平地面上，内壁为半圆柱面，截面如图所示，$A$~为半圆的最低点，$B$~为半圆水平直径的端点。凹槽恰好与竖直墙面接触，内有一质量为~$m$~的小滑块。用推力~$F$~推动小滑块由~$A$~点向~$B$~点缓慢移动，力~$F$~的方向始终沿圆弧的切线方向，在此过程中所有摩擦均可忽略，下列说法正确的是\gkblank

\gkfig{figures/q5-groove}

\begin{gkopts}(2)
  \task A．推力 $F$ 先增大后减小
  \task B．凹槽对滑块的支持力先减小后增大
  \task C．墙面对凹槽的压力先增大后减小
  \task D．水平地面对凹槽的支持力先减小后增大
\end{gkopts}
\end{question}

\begin{question}
\gkti{6}如图所示，理想变压器原、副线圈匝数比为~$n_{1}:n_{2}$，输入端~$C$、$D$~接入电压有效值恒定的交流电源，灯泡~$L_{1}$、$L_{2}$~的阻值始终与定值电阻~$R_{0}$~的阻值相同。在滑动变阻器~$R$~的滑片从~$a$~端滑动到~$b$~端的过程中，两个灯泡始终发光且工作在额定电压以内，下列说法正确的是\gkblank

\gkfig{figures/q6-transformer}

\begin{gkopts}(2)
  \task A．$L_{1}$ 先变暗后变亮，$L_{2}$ 一直变亮
  \task B．$L_{1}$ 先变亮后变暗，$L_{2}$ 一直变亮
  \task C．$L_{1}$ 先变暗后变亮，$L_{2}$ 先变亮后变暗
  \task D．$L_{1}$ 先变亮后变暗，$L_{2}$ 先变亮后变暗
\end{gkopts}
\end{question}

\textbf{（二）多项选择题（本题共 4 小题，每小题 5 分，共 20 分。在每小题给出的四个选项中，有多项符合题目要求。全部选对的得 5 分，选对但不全的得 3 分，有选错的得 0 分。）}

\begin{question}
\gkti{7}2021 年 4 月 29 日，中国空间站天和核心舱发射升空，准确进入预定轨道。根据任务安排，后续将发射问天实验舱和梦天实验舱，计划 2022 年完成空间站在轨建造。核心舱绕地球飞行的轨道可视为圆轨道，轨道离地面的高度约为地球半径的~$\dfrac{1}{16}$。下列说法正确的是\gkblank

\begin{gkopts}(1)
  \task A．核心舱进入轨道后所受地球的万有引力大小约为它在地面时的 $\left(\dfrac{16}{17}\right)^{2}$ 倍
  \task B．核心舱在轨道上飞行的速度大于 \SI{7.9}{\kilo\metre\per\second}
  \task C．核心舱在轨道上飞行的周期小于 \SI{24}{\hour}
  \task D．后续加挂实验舱后，空间站由于质量增大，轨道半径将变小
\end{gkopts}
\end{question}

\begin{question}
\gkti{8}如图甲所示，质量分别为~$m_{A}$、$m_{B}$~的~$A$、$B$~两物体用轻弹簧连接构成一个系统，外力~$F$~作用在~$A$~上，系统静止在光滑水平面上（$B$~靠墙面），此时弹簧形变量为~$x$。撤去外力并开始计时，$A$、$B$~两物体运动的~$a$--$t$~图像如图乙所示，$S_{1}$~表示 0 到~$t_{1}$~时间内~$A$~的~$a$--$t$~图线与坐标轴所围面积大小，$S_{2}$、$S_{3}$~分别表示~$t_{1}$~到~$t_{2}$~时间内~$A$、$B$~的~$a$--$t$~图线与坐标轴所围面积大小。$A$~在~$t_{1}$~时刻的速度为~$v_{0}$。下列说法正确的是\gkblank

\gkfig{figures/q8-spring}

甲

\gkfig{figures/q8-at}

乙

\begin{gkopts}(2)
  \task A．0 到 $t_{1}$ 时间内，墙对 $B$ 的冲量等于 $m_{A}v_{0}$
  \task B．$m_{A}>m_{B}$
  \task C．$B$ 运动后，弹簧的最大形变量等于 $x$
  \task D．$S_{1}-S_{2}=S_{3}$
\end{gkopts}
\end{question}

\begin{question}
\gkti{9}如图所示，圆心为~$O$~的圆处于匀强电场中，电场方向与圆平面平行，$ab$~和~$cd$~为该圆直径。将电荷量为~$q$（$q>0$）的粒子从~$a$~点移动到~$b$~点，电场力做功为~$2W$（$W>0$）；若将该粒子从~$c$~点移动到~$d$~点，电场力做功为~$W$。下列说法正确的是\gkblank

\gkfig{figures/q9-circle}

\begin{gkopts}(1)
  \task A．该匀强电场的电场强度方向与 $ab$ 平行
  \task B．将该粒子从 $d$ 点移动到 $b$ 点，电场力做功为 $0.5W$
  \task C．$a$ 点电势低于 $c$ 点电势
  \task D．若只受电场力，从 $d$ 点射入圆形电场区域的所有带电粒子都做曲线运动
\end{gkopts}
\end{question}

\begin{question}
\gkti{10}两个完全相同的正方形匀质金属框，边长为~$l$，通过长为~$l$~的绝缘轻质杆相连，构成如图所示的组合体。距离组合体下底边~$h$~处有一方向水平、垂直于纸面向里的匀强磁场。磁场区域上下边界水平，高度为~$l$，左右宽度足够大。把该组合体在垂直于磁场的平面内以初速度~$v_{0}$~水平无旋转抛出，设置合适的磁感应强度大小~$B$~使其匀速通过磁场，不计空气阻力。下列说法正确的是\gkblank

\gkfig{figures/q10-frame}

\begin{gkopts}(1)
  \task A．$B$ 与 $v_{0}$ 无关，与 $\sqrt{h}$ 成反比
  \task B．通过磁场的过程中，金属框中电流的大小和方向保持不变
  \task C．通过磁场的过程中，组合体克服安培力做功的功率与重力做功的功率相等
  \task D．调节 $h$、$v_{0}$ 和 $B$，只要组合体仍能匀速通过磁场，则其通过磁场的过程中产生的热量不变
\end{gkopts}
\end{question}

\gkpart{二、实验题}

\textbf{（一）必考题（本题共 2 小题，共 15 分。）}

\begin{question}
\gkti{11}某实验小组利用图甲所示装置探究加速度与物体所受合外力的关系。主要实验步骤如下：

\gkphoto[0.5\textwidth]{photos/image12-airtrack.png}{气垫导轨实验装置（原图保留，非矢量例外）}

甲

\begin{enumerate}[leftmargin=2em]
  \item 用游标卡尺测量垫块厚度~$h$，示数如图乙所示，$h=$\gkblank~cm。
  \item 接通气泵，将滑块轻放在气垫导轨上，调节导轨至水平。
  \item 在右支点下放一垫块，改变气垫导轨的倾斜角度。
  \item 在气垫导轨合适位置释放滑块，记录垫块个数~$n$~和滑块对应的加速度~$a$。
  \item 在右支点下增加垫块个数（垫块完全相同），重复步骤 4，记录数据如下表：
  \item 根据表中数据在图丙上描点，绘制图线。
\end{enumerate}

\gkphoto[0.45\textwidth]{photos/image13-vernier.png}{游标卡尺测量垫块厚度（原图保留，非矢量例外）}

乙

\begin{table}[ht]
\centering
\caption{垫块个数与加速度测量数据}
\small
\begin{tabular}{ccccccccc}
\toprule
$n$ & 1 & 2 & 3 & 4 & 5 & 6\\
\midrule
$a$/(m$\cdot$s$^{-2}$) & 0.087 & 0.180 & 0.260 &  & 0.425 & 0.519\\
\bottomrule
\end{tabular}
\end{table}

\gkfig{figures/q11-grid}

丙

如果表中缺少的第 4 组数据是正确的，其应该是\gkblank~m/s$^{2}$（保留三位有效数字）。
\end{question}

\begin{question}
\gkti{12}某实验小组需测定电池的电动势和内阻，器材有：一节待测电池、一个单刀双掷开关、一个定值电阻（阻值为~$R_{0}$）、一个电流表（内阻为~$R_{A}$）、一根均匀电阻丝（电阻丝总阻值大于~$R_{0}$，并配有可在电阻丝上移动的金属夹）、导线若干。由于缺少刻度尺，无法测量电阻丝长度，但发现桌上有一个圆形时钟表盘。某同学提出将电阻丝绕在该表盘上，利用圆心角来表示接入电路的电阻丝长度。主要实验步骤如下：

\gkphoto[0.5\textwidth]{photos/image16-wiring.png}{实验电路实物连接（原图保留，非矢量例外）}

甲

\begin{enumerate}[leftmargin=2em]
  \item 将器材如图甲连接。
  \item 开关闭合前，金属夹应夹在电阻丝的\gkblank（选填``$a$''或``$b$''）端。
  \item 改变金属夹的位置，闭合开关，记录每次接入电路的电阻丝对应的圆心角~$\theta$~和电流表示数~$I$，得到多组数据。
  \item 整理数据并在坐标纸上描点绘图，所得图像如图乙所示，图线斜率为~$k$，与纵轴截距为~$d$，设单位角度对应电阻丝的阻值为~$r_{0}$，该电池电动势和内阻可表示为 $E=$\gkblank，$r=$\gkblank（用 $R_{0}$、$R_{A}$、$k$、$d$、$r_{0}$ 表示）。
  \item 为进一步确定结果，还需要测量单位角度对应电阻丝的阻值~$r_{0}$。利用现有器材设计实验，在图丙方框中画出实验电路图（电阻丝用滑动变阻器符号表示）。
  \item 利用测出的 $r_{0}$，可得该电池的电动势和内阻。
\end{enumerate}

\gkfig{figures/q12-1overI}

乙

\gkfig{figures/q12-box}

丙
\end{question}

\gkpart{三、解答题}

\textbf{（一）必考题（本题共 2 小题，共 28 分。）}

\begin{question}
\gkti{13}带电粒子流的磁聚焦和磁控束是薄膜材料制备的关键技术之一。带电粒子流（每个粒子的质量为~$m$、电荷量为~$+q$）以初速度~$v$~垂直进入磁场，不计重力及带电粒子之间的相互作用。对处在~$xOy$~平面内的粒子，求解以下问题。

\gkfig{figures/q13-circ}

甲

\begin{enumerate}[leftmargin=2em, label=(\arabic*)]
  \item 如图甲所示，宽度为 $2r_{1}$ 的带电粒子流沿 $x$ 轴正方向射入圆心为 $A(0,r_{1})$、半径为 $r_{1}$ 的圆形匀强磁场中，若带电粒子流经过磁场后都汇聚到坐标原点 $O$，求该磁场磁感应强度 $B_{1}$ 的大小。
  \item 如图甲所示，虚线框为边长等于 $2r_{2}$ 的正方形，其几何中心位于 $C(0,-r_{2})$。在虚线框内设计一个区域面积最小的匀强磁场，使汇聚到 $O$ 点的带电粒子流经过该区域后宽度变为 $2r_{2}$，并沿 $x$ 轴正方向射出。求该磁场磁感应强度 $B_{2}$ 的大小和方向，以及该磁场区域的面积（无需写出面积最小的证明过程）。
\end{enumerate}

\gkfig{figures/q13-squares}

乙

\begin{enumerate}[leftmargin=2em, label=(\arabic*), resume]
  \item 如图乙所示，虚线框 Ⅰ 和 Ⅱ 均为边长等于 $r_{3}$ 的正方形，虚线框 Ⅲ 和 Ⅳ 均为边长等于 $r_{4}$ 的正方形。在 Ⅰ、Ⅱ、Ⅲ 和 Ⅳ 中分别设计一个区域面积最小的匀强磁场，使宽度为 $2r_{3}$ 的带电粒子流沿 $x$ 轴正方向射入 Ⅰ 和 Ⅱ 后汇聚到坐标原点 $O$，再经过 Ⅲ 和 Ⅳ 后宽度变为 $2r_{4}$，并沿 $x$ 轴正方向射出，从而实现带电粒子流的同轴控束。求 Ⅰ 和 Ⅲ 中磁场磁感应强度的大小，以及 Ⅱ 和 Ⅳ 中匀强磁场区域的面积（无需写出面积最小的证明过程）。
\end{enumerate}
\end{question}

\begin{question}
\gkti{14}如图所示，竖直平面内一足够长的光滑倾斜轨道与一长为~$l$~的水平轨道通过一小段光滑圆弧平滑连接，水平轨道右下方有一段弧形轨道~$PQ$。质量为~$m$~的小物块~$A$~与水平轨道间的动摩擦因数为~$\mu$。以水平轨道末端~$O$~点为坐标原点建立平面直角坐标系~$xOy$，$x$~轴的正方向水平向右，$y$~轴的正方向竖直向下，弧形轨道~$P$~端坐标为~$(2\mu l,\,\mu l)$，$Q$~端在~$y$~轴上。重力加速度为~$g$。

\gkfig{figures/q14-track}

\begin{enumerate}[leftmargin=2em, label=(\arabic*)]
  \item 若 $A$ 从倾斜轨道上距 $x$ 轴高度为 $2\mu l$ 的位置由静止开始下滑，求 $A$ 经过 $O$ 点时的速度大小。
  \item 若 $A$ 从倾斜轨道上不同位置由静止开始下滑，经过 $O$ 点落在弧形轨道 $PQ$ 上的动能均相同，求 $PQ$ 的曲线方程。
  \item 将质量为 $\lambda m$（$\lambda$ 为常数且 $\lambda\ge5$）的小物块 $B$ 置于 $O$ 点，$A$ 沿倾斜轨道由静止开始下滑，与 $B$ 发生弹性碰撞（碰撞时间极短），要使 $A$ 和 $B$ 均能落在弧形轨道上，且 $A$ 落在 $B$ 落点的右侧，求 $A$ 下滑的初始位置距 $x$ 轴高度的取值范围。
\end{enumerate}
\end{question}

\textbf{（二）选考题（本题共 1 小题，共 13 分。请考生从两道选考题中任选一题作答，本卷转录第 15 题。）}

\begin{question}
\gkti{15}[物理——选修 3-3]

(1)（5 分）如图所示，两端开口、下端连通的导热汽缸，用两个轻质绝热活塞（截面积分别为~$S_{1}$~和~$S_{2}$）封闭一定质量的理想气体，活塞与汽缸壁间无摩擦。在左端活塞上缓慢加细沙，活塞从~$A$~下降~$h$~高度到~$B$~位置时，活塞上细沙的总质量为~$m$。在此过程中，用外力~$F$~作用在右端活塞上，使活塞位置始终不变。整个过程环境温度和大气压强~$p_{0}$~保持不变，系统始终处于平衡状态，重力加速度为~$g$。下列说法正确的是\gkblank（填正确答案标号。选对 1 个得 2 分，选对 2 个得 4 分，选对 3 个得 5 分。每选错 1 个扣 3 分，最低得分为 0 分）

\gkfig{figures/q15-cylinder}

\begin{gkopts}(1)
  \task A．整个过程，外力 $F$ 做功大于 0，小于 $mgh$
  \task B．整个过程，理想气体的分子平均动能保持不变
  \task C．整个过程，理想气体的内能增大
  \task D．整个过程，理想气体向外界释放的热量小于 $p_{0}S_{1}h+mgh$
  \task E．左端活塞到达 $B$ 位置时，外力 $F$ 等于 $\dfrac{mgS_{2}}{S_{1}}$
\end{gkopts}

(2)（8 分）小赞同学设计了一个用电子天平测量环境温度的实验装置，如图所示。导热汽缸开口向上并固定在桌面上，用质量~$m_{1}=\SI{600}{\gram}$、截面积~$S=\SI{20}{\centi\metre\squared}$~的活塞封闭一定质量的理想气体，活塞与汽缸壁间无摩擦。一轻质直杆中心置于固定支点~$A$~上，左端用不可伸长的细绳竖直悬挂活塞，右端用相同细绳竖直悬挂一个质量~$m_{2}=\SI{1200}{\gram}$~的铁块，并将铁块放置到电子天平上。当电子天平示数为 \SI{600.0}{\gram} 时，测得环境温度~$T_{1}=\SI{300}{\kelvin}$。设外界大气压强~$p_{0}=\SI{1.0e5}{\pascal}$，重力加速度~$g$~取 \SI{10}{\metre\per\second\squared}。

\gkfig{figures/q15-balance}

\begin{enumerate}[leftmargin=2em, label=(\roman*)]
  \item 当电子天平示数为 \SI{400.0}{\gram} 时，环境温度 $T_{2}$ 为多少？
  \item 该装置可测量的最高环境温度 $T_{\max}$ 为多少？
\end{enumerate}
\end{question}
